\documentclass[11pt]{article}
\usepackage[a4paper,margin=24mm]{geometry}
\usepackage{amsmath,amssymb,xcolor}
\usepackage[colorlinks=true,linkcolor=blue!45!black,urlcolor=blue!45!black]{hyperref}
\setlength{\parindent}{0pt}
\setlength{\parskip}{6pt}
\setlength{\emergencystretch}{2em}
\newcommand{\Q}{\mathbb Q}
\newcommand{\Z}{\mathbb Z}
\newcommand{\R}{\mathbb R}
\newcommand{\dd}{\,d}
\newcommand{\Hfive}[1]{H_{#1}^{(5)}}
\newcommand{\source}[1]{{\small\color{blue!45!black}#1}\par}
\pagestyle{plain}
\begin{document}
\begin{center}
{\LARGE Why the determinant construction proves\\[3pt]
$\zeta(5)$ irrational}\\[9pt]
{\large An overview following Chun Wing Lai's exposition}\\[5pt]
Draft \enspace$\cdot$\enspace 14 November 2026
\end{center}

The argument constructs integer polynomials whose values at $\zeta(5)$
become extraordinarily small. Their degrees grow only linearly, while their
values shrink exponentially in the square of the size parameter. If
$\zeta(5)$ were rational, clearing its fixed denominator would then produce
a positive integer tending to zero.

We explain the construction and this final contradiction, taking two precise
estimates as inputs. The analytic estimate is Lai's Proposition 11.7.
The arithmetic estimate is Fauzan's original bound, assumed separately:
it does not follow from Lai's sections 8--11 as literally printed.
The source note at the end explains this qualification. Both estimates
are stated before being used.

\section*{1. What would be enough?}
\source{Lai, section 3, Lemma 3.1.}
Write $\xi=\zeta(5)$. Suppose that, for every sufficiently large positive
integer $n$, we can find a polynomial $Q_n\in\Z[X]$ satisfying
\[
 \deg Q_n\leq37n,\qquad 0<Q_n(\xi)\leq e^{-cn^2},
 \qquad c>0\text{ fixed}.
\]
The number $37$ will come from the size of a matrix. Its particular value
does not matter to the contradiction; what matters is linear growth.

Assume for a moment that $\xi=a/b$, with $a\in\Z$ and $b\in\Z_{>0}$.
If $Q_n(X)=\sum_{j=0}^{37n}q_jX^j$, then
\[
 b^{37n}Q_n(a/b)=\sum_{j=0}^{37n}q_j a^j b^{37n-j}\in\Z.
\]
This integer is positive, but the small-value estimate gives
\[
 0<b^{37n}Q_n(a/b)
 \leq \exp\bigl(37n\log b-cn^2\bigr)\longrightarrow0.
\]
Eventually it would lie strictly between zero and one. That is impossible.

The construction must therefore accomplish two things simultaneously.
It must produce a polynomial with integer coefficients, and it must keep
the value of that polynomial at $\xi$ small without letting it vanish.
These demands pull in different directions: multiplying by a large integer
clears denominators, but also enlarges the value we are trying to make small.

A determinant gives enough room to satisfy both demands. Its entries will
be linear expressions in $X$ with rational coefficients. At $X=\xi$, the
same entries become moments of a positive measure. The first description
provides arithmetic structure; the second makes nonvanishing visible and
opens the way to an analytic estimate. The rest of the proof is the passage
between these two descriptions.

\newpage
\section*{2. One rule connects rational coefficients and a positive integral}
\source{Lai, section 4, Definition 4.2 and Lemmas 4.5--4.6.}
For $j\geq1$, let
\[
 \Hfive{j}=\sum_{v=1}^j\frac1{v^5},\qquad \Hfive{0}=0.
\]
Consider rational functions $R\in\Q(t)$ with simple poles only at
$-1^2,-2^2,-3^2,\ldots$. Partial fractions write each one uniquely as
\[
 R(t)=P(t)+\sum_j\frac{c_j}{t+j^2},
\]
where $P$ is a polynomial, all coefficients are rational, and the sum is
finite. A linear rule on these functions is determined once we specify it
on powers of $t$ and on the simple fractions.

Use Bernoulli numbers with $B_1=-1/2$, and define
\begin{align*}
 \mu_X(t^e)&=\frac{(-1)^eB_{2e+2}(2e+3)(2e+4)(2e+5)}{24},
       && e\geq0,\\
 \mu_X\!\left(\frac1{t+j^2}\right)
   &=j^4\bigl(X-\Hfive{j}\bigr)-\frac14+\frac1{2j},
       && j\geq1.
\end{align*}
Extend by linearity. Every value $\mu_X(R)$ is an affine polynomial in $X$
with rational coefficients. The polynomial part of $R$ contributes no $X$;
all dependence on $X$ comes from its poles. For example,
\[
 \mu_X(1)=\frac5{12},\qquad
 \mu_X\!\left(\frac1{t+1}\right)=X-\frac34.
\]

The reason for these formulas is that at $X=\xi$ they describe integration:
\[
 \boxed{\mu_\xi(R)=\int_0^\infty R(y^2)w(y)\dd y},\qquad
 w(y)=\frac{(2\pi)^4y^5}{12}
       \sum_{\ell\geq1}\ell^4e^{-2\pi\ell y}>0.
\]
The weight is bounded near zero and decays exponentially at infinity, so
these integrals converge. To check the identity, it is enough to check the
same two types of building block.

For powers of $t$, integrate each exponential and then use Euler's formula
for even zeta values:
\[
 \int_0^\infty y^{2e}w(y)\dd y
 =\frac{(2e+5)!}{12(2\pi)^{2e+2}}\zeta(2e+2)
 =\mu_X(t^e).
\]
For a simple pole, the Hurwitz zeta integral gives
\[
 \int_0^\infty\frac{w(y)}{y^2+j^2}\dd y
 =j^4\zeta(5,j)-\frac1{2j}-\frac14
 =j^4\bigl(\xi-\Hfive{j}\bigr)-\frac14+\frac1{2j}.
\]
Here $\zeta(5,j)=\sum_{v=0}^\infty(j+v)^{-5}
=\xi-\Hfive{j}+j^{-5}$. Linearity completes the check.

Thus $X$ has a specific role: it stands in for the one zeta value occurring
in the pole integrals. Keeping $X$ symbolic gives rational algebra;
putting $X=\xi$ gives a positive integral. We will use both descriptions
of the same rule.

\newpage
\section*{3. The determinant is a polynomial, and its value cannot vanish}
\source{Lai, sections 5--6; especially Corollary 5.5 and Proposition 5.6.}
For a positive integer $n$, set
\[
 K=40n,\qquad N=3n,\qquad h=K-N=37n,
 \qquad D_m(t)=\prod_{j=1}^m(t+j^2).
\]
Introduce the rational multiplier
\[
 A_K(t)=\frac{D_N(t)^6}{D_K(t)}
       =\frac{D_N(t)^5}{\displaystyle\prod_{k=N+1}^K(t+k^2)}.
\]
The cancellation shows that its remaining poles are simple. Also,
$A_K(y^2)>0$ for $y>0$. These two observations let us apply both the
algebraic and the integral descriptions of $\mu$.

Build the $h$ by $h$ matrix and its determinant:
\[
 G_K(X)=\bigl[\mu_X(A_K(t)t^{i+j})\bigr]_{0\leq i,j<h},
 \qquad \Delta_K(X)=\det G_K(X).
\]
The entries depend only on $i+j$, which is why the matrix is called
\emph{Hankel}. Each entry is affine in $X$. Every term in the determinant
is a product of $h$ such entries, so
\[
 \Delta_K\in\Q[X],\qquad \deg\Delta_K\leq h=37n.
\]
This elementary degree bound is all the irrationality argument needs.
Lai proves the stronger equality using a Vandermonde determinant on the
$h$ remaining pole locations.

Now evaluate at $X=\xi$. The entries become
\[
 G_K(\xi)_{ij}=\int_0^\infty
        y^{2i}y^{2j}A_K(y^2)w(y)\dd y.
\]
This is a \emph{Gram matrix}: it records inner products of
$1,y^2,\ldots,y^{2h-2}$ for a positive weight. More explicitly, if
$q(t)=\sum_{i=0}^{h-1}v_it^i$ is a nonzero real polynomial, then
\[
 v^{\mathsf T}G_K(\xi)v
 =\int_0^\infty q(y^2)^2A_K(y^2)w(y)\dd y>0.
\]
The inequality is strict because the weight is positive and a nonzero
polynomial has only finitely many zeros. Thus $G_K(\xi)$ is positive
definite, and $\Delta_K(\xi)>0$. This is the key safeguard against a
determinant that is small merely because it vanishes.

The estimates we will borrow are written for an explicitly rescaled
determinant:
\[
 F_K(X)=S_K\Delta_K(X),\qquad
 S_K=\frac{(K!)^{2h}4^{h-1}}
 {(N!)^{12h}\displaystyle\prod_{i=1}^{h-1}((2i)!)^2}>0.
\]
Since $S_K$ is a positive rational number independent of $X$, we still have
$F_K\in\Q[X]$, $\deg F_K\leq37n$, and $F_K(\xi)>0$.
The factorial expression puts the arithmetic and analytic estimates on a
common scale. It does not, by itself, make every coefficient integral.

\newpage
\section*{4. The two estimates we take as inputs}
We now state the quantitative facts supplied by the arithmetic and analytic
parts of the argument. We use slightly weaker, rounded constants so that
the final comparison is transparent. Throughout, $K$ tends to infinity
through positive multiples of $40$.

\textbf{Arithmetic estimate.}
For all sufficiently large $K$, there is a positive rational number $m_K$
such that
\[
 \boxed{m_KF_K\in\Z[X],\qquad \log m_K\leq1.35K^2.}
\]
This is the normalization estimate underlying Lai's sections 8--10,
with the cutoff fixed at $M=200$. We take its precise arithmetic form
from Fauzan's Proposition 5.1 and equation (5.21); the source note explains
why this specification matters.

The content is the bound on the cost, not just the existence of a common
denominator. For each prime, the arithmetic argument bounds its exponent
in every coefficient. Negative exponents require clearing denominators;
positive exponents allow common factors to be divided out. This explains
why $m_K$ is positive and rational, but need not be an integer. Summing the
local contributions controls $\log m_K$.

Section 7 prepares that calculation by rewriting the same functional as
$\mu_X(R)=\tau_X(x^5R(-x^2))$. The substitution replaces poles at negative
squares by poles at integers. It makes residue classes modulo a prime
available for the divisibility argument. For the present overview, all
the resulting arithmetic is contained in the estimate just stated.

\textbf{Analytic estimate.}
For all sufficiently large $K$,
\[
 \boxed{0<F_K(\xi)\leq\exp(-1.36K^2).}
\]
This is a weakened form of Lai's Proposition 11.7. Its stronger statement is
\[
 \log F_K(\xi)\leq
 -1.3669955K^2+24K\log K+200K.
\]
After division by $K^2$, the last two terms tend to zero. The strict gap
between $1.3669955$ and $1.36$ therefore gives the displayed estimate.

One identity explains why the determinant is suited to this scale of decay.
Andr\'eief's identity, in the form needed here, states
\[
 \Delta_K(\xi)=\frac1{h!}\int_{(0,\infty)^h}
 \prod_{i<j}(y_i^2-y_j^2)^2
 \prod_{i=1}^h A_K(y_i^2)w(y_i)\dd y_1\cdots\dd y_h.
\]
Taking logarithms turns the pairwise factors into a sum of logarithmic
distances. There are on the order of $h^2$ such pairs. The remaining
weights contribute a competing penalty. Section 11 bounds this balance
by a logarithmic-energy argument and combines it with the factorial
rescaling $S_K$. The analytic estimate above is the output we need;
its proof is being taken as given.

\newpage
\section*{5. The arithmetic cost is smaller than the analytic decay}
Define
\[
 Q_n(X)=m_{40n}F_{40n}(X).
\]
The arithmetic estimate gives $Q_n\in\Z[X]$. The positive scalar $m_{40n}$
preserves both the degree bound and positivity, so
\[
 \deg Q_n\leq37n,\qquad Q_n(\xi)>0.
\]
For sufficiently large $n$, put $K=40n$ and add the two logarithmic estimates:
\begin{align*}
 \log Q_n(\xi)
 &=\log m_K+\log F_K(\xi)\\
 &\leq1.35K^2-1.36K^2
 =-\frac{K^2}{100}
 =-16n^2.
\end{align*}
Consequently
\[
 \boxed{Q_n\in\Z[X],\qquad \deg Q_n\leq37n,
        \qquad 0<Q_n(\xi)\leq e^{-16n^2}.}
\]
These are exactly the properties sought at the beginning. If $\xi=a/b$,
then $b^{37n}Q_n(a/b)$ is a positive integer and
\[
 0<b^{37n}Q_n(a/b)
 \leq\exp\bigl(37n\log b-16n^2\bigr)\longrightarrow0.
\]
The contradiction proves that $\zeta(5)$ is irrational, subject to the two
stated estimates.

The construction works because the determinant has modest algebraic degree
but enough analytic decay to survive integer normalization. The positive
moment representation prevents vanishing; the arithmetic estimate keeps
the normalization affordable. The numerical margin need only be strictly
positive. There is no need to retain the sharper decay constant, or to
control the polynomial near $\xi$, for this irrationality conclusion.

\bigskip
\textbf{Source note: the arithmetic input.}
This overview follows Chun Wing Lai,
\href{https://ai.vixra.org/pdf/2609.0080v1.pdf}{\emph{On Fauzan's Proof that
$\zeta(5)$ is Irrational}}, v1, 29 September 2026.
The arithmetic input used above is the statement
\[
 \limsup_{\substack{K\to\infty\\40\mid K}}
       \frac{\log m_{K,200}}{K^2}
 \leq A_{200}=1.3495876977\ldots<1.35,
\]
from Aabir Fauzan's
\href{https://zenodo.org/records/22826419}{17 September 2026 manuscript},
equation (5.21), with the constant evaluated in Appendix B. The strict
inequality implies the eventual bound used on the preceding page.

Lai's Theorem 9.3 and equations (106)--(107) instead include an additional
$+E_M$. His following prose treats that positive term as a subtraction.
The two formulations cannot be used interchangeably. Keeping the printed
correction would require a further bound
$E_{200}<0.017407802258\ldots$ to obtain a negative total exponent and hence
some decay constant $c>0$, not necessarily $16$.
The original arithmetic estimate is an additional assumed input; citing it
does not verify or repair the disputed arithmetic. With Lai's literal
corrected statement alone, the final numerical comparison remains
conditional on the additional bound.
\end{document}
